\section{The Collision That Is Still There}
\label{sec:ch09-the-collision-that-is-still-there}

\index{seam!the collision chapter 7 left|(}%
Chapter~\ref{ch:composition} finished with one error it could not fix.
Compose the Sessions subgraph against the hand-written Speakers document again
and it is where it was left:

\begin{minted}[fontsize=\footnotesize,breakanywhere]{text}
The Object "Speaker" defines the same fields in multiple subgraphs without the "@shareable" directive:
 The field "name" is defined in the following subgraphs: "sessions", "speakers".
 However, it is not declared "@shareable" in any of them.
 The field "bio" is defined in the following subgraphs: "sessions", "speakers".
 However, it is not declared "@shareable" in any of them.
\end{minted}

\index{seam!has two fixes and only one is right}%
There are two ways to make that message go away and only one of them is
federation. Declaring both fields \code{@shareable} tells the composer that
either service may answer for a speaker's name, which is a promise neither can
keep and the exact thing chapter~\ref{ch:composition} closed by arguing
against. The other way is to make it true that one service owns the rows.

\index{ownership!is a decision the composer cannot make}%
The composer cannot help you choose. It can see that two documents claim one
field; which of the two should give it up is a question about who owns the
data, and chapter~\ref{ch:ownership} is a whole chapter on how badly that
question goes when nobody answers it deliberately. Here the answer is easy,
because the conference has been running the other way round since
chapter~\ref{ch:hot-chocolate-zero}: the Sessions service stores speakers only
because it was the first service there was.

\subsection{What Moves, and What Cannot}

\index{seam!the row moves and the key stays}%
The rows move. The \code{SpeakerId} column on each session does not, and could
not: it is how a session says which speaker gives it, and a session is
Sessions' to describe. So the split runs between the identifier and everything
the identifier points at. Sessions keeps a number. Speakers gets the name, the
biography and the table.

\index{seam!the field survives the split}%
The field a client asks for survives that intact. \code{Session.speaker} still
exists, still returns a \code{Speaker}, and still answers. What changes is that
the object it returns is assembled by two processes, and that the router's
query plan says so out loud.

\index{ownership!two subgraphs share no code}%
One rule governs everything that follows, and it comes before the first file
because everything after it depends on the rule: the two services share no code.
Not a model library, not a \code{Common} project, not a linked file. Where
both need the same shape, each declares its own copy. These would be separate
repositories owned by separate teams in any conference big enough to need
this, and a shared type would make the book demonstrate the opposite of what
it says. It also costs something real, and the next section is where the
bill arrives: a DataLoader, an interceptor and a statement log all get typed
out a second time.
\index{seam!the collision chapter 7 left|)}%
